%%%%%%%%%%%%%%%%%%%%%%%%%%%%%%%%%%%%%%%%%%%%%%%%%%%%%%%%%%%%%%%%%%%%%%
% Lớp 9 | Chương V. Đường tròn | Vị trí tương đối của đường thẳng và đường tròn
% Dạng 9H5?16-3: Tính chất hai tiếp tuyến cắt nhau
%%%%%%%%%%%%%%%%%%%%%%%%%%%%%%%%%%%%%%%%%%%%%%%%%%%%%%%%%%%%%%%%%%%%%%


% nguồn: 2026-11-21_lop9_chuong5_bai14_17 (nhập 21/11/2026)
\begin{ex}%[9H5N16-3]%[ma:faac3272]
% nguon: SGK Toán 9 KNTT tập 1 – Bài 16, Vị trí tương đối của đường thẳng và đường tròn
	Từ điểm $M$ nằm ngoài đường tròn $(O)$, kẻ hai tiếp tuyến $MA, MB$ đến đường tròn ($A, B$ là các tiếp điểm). Khẳng định nào sau đây \textbf{sai}?
	\choice
	{$MA = MB$}
	{$MO$ là tia phân giác của $\widehat{AMB}$}
	{$OM$ là tia phân giác của $\widehat{AOB}$}
	{\True $OA = MA$}
	\loigiai{
		Theo tính chất hai tiếp tuyến cắt nhau, ta có $MA = MB$, $MO$ là tia phân giác của $\widehat{AMB}$ và $OM$ là tia phân giác của $\widehat{AOB}$. Độ dài $OA$ và $MA$ không nhất thiết bằng nhau (chỉ bằng nhau khi $\widehat{AMB} = 90^\circ$).
	}
\end{ex}

% nguồn: 2026-11-21_lop9_chuong5_bai14_17 (nhập 21/11/2026)
\begin{ex}%[9H5H16-3]%[ma:1754f79e]
% nguon: SGK Toán 9 KNTT tập 1 – Bài 16, Vị trí tương đối của đường thẳng và đường tròn
	Từ điểm $M$ nằm ngoài đường tròn $(O; 3\text{ cm})$ sao cho $OM = 6\text{ cm}$, kẻ hai tiếp tuyến $MA, MB$ đến đường tròn ($A, B$ là các tiếp điểm). Xét tính đúng, sai của các khẳng định sau:
	\choiceTF
	{\True $MA = MB = 3\sqrt{3}\text{ cm}$}
	{\True $\widehat{AMO} = 30^\circ$}
	{Tam giác $MAB$ là tam giác vuông cân}
	{\True Chu vi tứ giác $OAMB$ bằng $(6 + 6\sqrt{3})\text{ cm}$}
	\loigiai{
		\begin{itemchoice}
			\itemch Đúng. Trong tam giác $OAM$ vuông tại $A$: $MA = \sqrt{OM^2 - OA^2} = \sqrt{6^2 - 3^2} = 3\sqrt{3}\text{ cm}$, và theo tính chất hai tiếp tuyến cắt nhau thì $MB = MA = 3\sqrt{3}\text{ cm}$.
			\itemch Đúng. Ta có $\sin \widehat{AMO} = \dfrac{OA}{OM} = \dfrac{3}{6} = \dfrac{1}{2}$ nên $\widehat{AMO} = 30^\circ$.
			\itemch Sai. Vì $MO$ là tia phân giác $\widehat{AMB}$ nên $\widehat{AMB} = 2 \cdot 30^\circ = 60^\circ$, kết hợp $MA = MB$ suy ra tam giác $MAB$ là tam giác đều.
			\itemch Đúng. Chu vi tứ giác $OAMB$ là $OA + OB + MA + MB = 3 + 3 + 3\sqrt{3} + 3\sqrt{3} = 6 + 6\sqrt{3}\text{ cm}$.
		\end{itemchoice}
	}
\end{ex}

% nguồn: 2026-11-21_lop9_chuong5_bai14_17 (nhập 21/11/2026)
\begin{bt}%[9H5V16-3]%[ma:c3d41774]
% nguon: SGK Toán 9 KNTT tập 1 – Bài 16, Vị trí tương đối của đường thẳng và đường tròn
	Từ điểm $M$ nằm ngoài đường tròn $(O; 6\text{ cm})$ sao cho $OM = 10\text{ cm}$, kẻ hai tiếp tuyến $MA, MB$ đến đường tròn $(O)$ với $A, B$ là các tiếp điểm. Tính độ dài dây $AB$ (đơn vị cm, viết kết quả dưới dạng số thập phân).
	\shortans{9,6}
	\loigiai{
		Trong tam giác $OAM$ vuông tại $A$: $MA = \sqrt{OM^2 - OA^2} = \sqrt{10^2 - 6^2} = 8\text{ cm}$.\\
		Vì $OA = OB$ và $MA = MB$ nên $OM$ là đường trung trực của đoạn $AB$. Gọi $H$ là giao điểm của $OM$ và $AB$, ta có $AH \perp OM$ và $AB = 2AH$.\\
		Trong tam giác $OAM$ vuông tại $A$ có đường cao $AH$:
		$$AH = \dfrac{OA \cdot MA}{OM} = \dfrac{6 \cdot 8}{10} = 4{,}8\text{ cm}.$$
		Suy ra $AB = 2 \cdot 4{,}8 = 9{,}6\text{ cm}$.
	}
\end{bt}

% nguồn: 2026-11-21_lop9_chuong5_bai16_17_bs (nhập 21/11/2026)
\begin{ex}%[9H5N16-3]%[ma:7d13ca2f]
% nguon: SGK Toán 9 KNTT tập 1 – Bài 16, Vị trí tương đối của đường thẳng và đường tròn
	Từ điểm $A$ nằm ngoài đường tròn $(O)$, kẻ hai tiếp tuyến $AB, AC$ đến đường tròn ($B, C$ là các tiếp điểm). Nếu $AB = 7\text{ cm}$ thì độ dài $AC$ bằng
	\choice
	{$3{,}5\text{ cm}$}
	{\True $7\text{ cm}$}
	{$14\text{ cm}$}
	{$7\sqrt{2}\text{ cm}$}
	\loigiai{
		Theo tính chất hai tiếp tuyến cắt nhau, điểm $A$ cách đều hai tiếp điểm $B$ và $C$, nên $AC = AB = 7\text{ cm}$.
	}
\end{ex}

% nguồn: 2026-11-21_lop9_chuong5_bai16_17_bs (nhập 21/11/2026)
\begin{ex}%[9H5N16-3]%[ma:04ca7167]
% nguon: SGK Toán 9 KNTT tập 1 – Bài 16, Vị trí tương đối của đường thẳng và đường tròn
	Từ điểm $M$ nằm ngoài đường tròn $(O)$, kẻ hai tiếp tuyến $MA, MB$ đến đường tròn ($A, B$ là các tiếp điểm). Biết $\widehat{AMO} = 25^\circ$, số đo góc $\widehat{AMB}$ bằng
	\choice
	{$25^\circ$}
	{\True $50^\circ$}
	{$65^\circ$}
	{$130^\circ$}
	\loigiai{
		Theo tính chất hai tiếp tuyến cắt nhau, $MO$ là tia phân giác của $\widehat{AMB}$, suy ra $\widehat{AMB} = 2\widehat{AMO} = 2 \cdot 25^\circ = 50^\circ$.
	}
\end{ex}

% nguồn: 2026-11-21_lop9_chuong5_bai16_17_bs (nhập 21/11/2026)
\begin{ex}%[9H5H16-3]%[ma:c16e2af4]
% nguon: SGK Toán 9 KNTT tập 1 – Bài 16, Vị trí tương đối của đường thẳng và đường tròn
	Từ điểm $A$ nằm ngoài đường tròn $(O; 5\text{ cm})$ sao cho $OA = 13\text{ cm}$, kẻ hai tiếp tuyến $AB, AC$ đến $(O)$ với $B, C$ là các tiếp điểm. Xét tính đúng, sai của các khẳng định sau:
	\choiceTF
	{\True $AB = AC = 12\text{ cm}$}
	{\True $OA$ là đường trung trực của đoạn thẳng $BC$}
	{Diện tích tứ giác $OBAC$ bằng $120\text{ cm}^2$}
	{\True Chu vi tứ giác $OBAC$ bằng $34\text{ cm}$}
	\loigiai{
		\begin{itemchoice}
			\itemch Đúng. Trong tam giác $OBA$ vuông tại $B$: $AB = \sqrt{13^2 - 5^2} = 12\text{ cm}$, suy ra $AC = AB = 12\text{ cm}$.
			\itemch Đúng. Vì $OB = OC$ và $AB = AC$ nên $OA$ là đường trung trực của $BC$.
			\itemch Sai. Diện tích tứ giác $OBAC$ bằng $2S_{OBA} = 2 \cdot \dfrac{1}{2} \cdot 5 \cdot 12 = 60\text{ cm}^2$.
			\itemch Đúng. Chu vi tứ giác $OBAC$ là $OB + OC + AB + AC = 5 + 5 + 12 + 12 = 34\text{ cm}$.
		\end{itemchoice}
	}
\end{ex}

% nguồn: 2026-11-21_lop9_chuong5_bai16_17_bs (nhập 21/11/2026)
\begin{ex}%[9H5H16-3]%[ma:87d5bdae]
% nguon: SGK Toán 9 KNTT tập 1 – Bài 16, Vị trí tương đối của đường thẳng và đường tròn
	Từ điểm $M$ nằm ngoài đường tròn $(O; 4\text{ cm})$, kẻ hai tiếp tuyến $MA, MB$ vuông góc với nhau ($A, B$ là các tiếp điểm). Xét tính đúng, sai của các khẳng định sau:
	\choiceTF
	{\True Tứ giác $OAMB$ là hình vuông}
	{\True Độ dài $MA = MB = 4\text{ cm}$}
	{Độ dài đoạn thẳng $OM$ bằng $8\text{ cm}$}
	{\True Diện tích phần nằm ngoài đường tròn $(O)$ nhưng bên trong tứ giác $OAMB$ (lấy $\pi \approx 3{,}14$, làm tròn đến hàng phần trăm) bằng $3{,}44\text{ cm}^2$}
	\loigiai{
		\begin{itemchoice}
			\itemch Đúng. Tứ giác $OAMB$ có $\widehat{A} = \widehat{B} = \widehat{M} = 90^\circ$ và $OA = OB = 4\text{ cm}$ nên là hình vuông cạnh $4\text{ cm}$.
			\itemch Đúng. Vì $OAMB$ là hình vuông cạnh $4\text{ cm}$ nên $MA = MB = 4\text{ cm}$.
			\itemch Sai. Đường chéo hình vuông cạnh $4\text{ cm}$ là $OM = 4\sqrt{2}\text{ cm}$.
			\itemch Đúng. Phần diện tích cần tìm bằng hiệu diện tích hình vuông $OAMB$ và hình quạt tròn $OAB$ góc $90^\circ$: $4^2 - \dfrac{90}{360}\pi \cdot 4^2 = 16 - 4\pi \approx 16 - 4 \cdot 3{,}14 = 3{,}44\text{ cm}^2$.
		\end{itemchoice}
	}
\end{ex}

% nguồn: 2026-11-21_lop9_chuong5_bai16_17_bs (nhập 21/11/2026)
\begin{bt}%[9H5V16-3]%[ma:a790ff3d]
% nguon: SGK Toán 9 KNTT tập 1 – Bài 16, Vị trí tương đối của đường thẳng và đường tròn
	Từ điểm $M$ nằm ngoài đường tròn $(O; 3\text{ cm})$ sao cho $OM = 5\text{ cm}$, kẻ hai tiếp tuyến $MA, MB$ đến đường tròn $(O)$ ($A, B$ là các tiếp điểm). Tính diện tích tứ giác $OAMB$ (đơn vị $\text{cm}^2$).
	\shortans{12}
	\loigiai{
		Trong tam giác $OAM$ vuông tại $A$: $MA = \sqrt{5^2 - 3^2} = 4\text{ cm}$.\\
		Vì $\triangle OAM = \triangle OBM$ nên diện tích tứ giác $OAMB$ là:
		$$S_{OAMB} = 2S_{OAM} = 2 \cdot \dfrac{1}{2} \cdot OA \cdot MA = 3 \cdot 4 = 12\text{ cm}^2.$$
	}
\end{bt}

% nguồn: 2026-11-21_lop9_chuong5_bai16_17_bs (nhập 21/11/2026)
\begin{bt}%[9H5V16-3]%[ma:b5ade249]
% nguon: SGK Toán 9 KNTT tập 1 – Bài 16, Vị trí tương đối của đường thẳng và đường tròn
	Cho đường tròn $(O)$ nội tiếp tam giác $ABC$, tiếp xúc với các cạnh $BC, CA, AB$ lần lượt tại các điểm $D, E, F$. Biết $AB = 7\text{ cm}$, $BC = 9\text{ cm}$, $CA = 8\text{ cm}$. Tính độ dài đoạn thẳng $AE$ (đơn vị cm).
	\shortans{3}
	\loigiai{
		Theo tính chất hai tiếp tuyến cắt nhau, đặt $AE = AF = x$, $BF = BD = y$, $CD = CE = z$.\\
		Ta có $x + y = AB = 7$, $y + z = BC = 9$, $z + x = CA = 8$.\\
		Cộng vế theo vế được $2(x + y + z) = 24 \Rightarrow x + y + z = 12$.\\
		Suy ra $AE = x = 12 - (y + z) = 12 - 9 = 3\text{ cm}$.
	}
\end{bt}
